\section{Discussion}
\label{sec:discussion}

Closure is a premise of the representation, not a result. Its scientific value
depends on what graph-valued composition enables beyond independent feature
extraction. An operator may expose a new structure, aggregate one exposed
earlier, or replace one view with another; it need not progressively refine a
representation. Additive union should agree with lossless feature concatenation
and serves as a parity check. Ordered and relational programs may expose
distinctions absent from that union because they construct new objects or
relations. An equally informed flat materialization determines whether those
distinctions require an interpretation graph or merely the same derived
information in another data structure. Empirical partitions and factor-specific
Pareto frontiers make these consequences observable without turning them into
algebraic axioms.

Mapped structure matters only if it improves diagnosis. We therefore require a
mapped witness for each representative distinction, test its overlap with the
known intervention, and compare it with provenance retained by the flat
control. When two signatures differ, the responsible
interpretation nodes can then be traced to concrete base subgraphs. Conversely,
identity width can collapse distinct certificates even when an operator has
retained them. Separating operator, certificate, hash, aggregation, and model
effects remains essential.

The minimal decomposition interface is an implementation advantage, not one of
the paper's evaluated scientific claims. Domain experts can contribute a
structural grouping without implementing the surrounding representation
machinery, but ease of authoring and automatically generated operators require
separate studies.

Threats to validity include the selected graph families, expression-selection
bias, canonicalization and attributed equivalence, imperfect baseline
reproduction, duplicate and multiplicity semantics, probe sensitivity, and
hardware-dependent costs. Synthetic controls establish mechanism but not
real-world usefulness. ZINC tests natural structural frequencies and scaling,
but it does not identify a causal structural distinction. Competitive
predictive learning and learned explanations remain subjects of later work.
